\documentclass[10pt,a4paper]{amsart}
\usepackage[margin=19mm]{geometry}
\usepackage{amsmath,amssymb,mathtools}
\usepackage[scr=rsfs,cal=euler]{mathalfa}
\usepackage{xfrac}
\usepackage[unicode,hidelinks,bookmarks=false]{hyperref}
\newcommand{\ie}{i.e.\ }
\newcommand{\cf}{cf.\ }
\newcommand{\resp}{respectively\ }
\newcommand{\Subsection}{\subsection}
\providecommand{\nobreakdash}{\nobreak\hbox{-}}
\setlength{\emergencystretch}{2em}
\multlinegap0pt
\usepackage{bm}
\usepackage{bbm}
\usepackage{dsfont}
\usepackage{xspace}
\usepackage{cancel}
\usepackage{mathtools}
\usepackage{amsthm}
\makeatletter
\@ifundefined{claim}{\newtheorem{claim}{Claim}}{}
\makeatother
\newcommand{\then}{\mathbin{\supset}}
\title[Incompleteness in Quantified Conditional Logic]{Incompleteness in Quantified Conditional Logic}
\author{Alexander W. Kocurek, James Walsh, Yale Weiss}
\date{Source: arXiv:2602.04073v1 (CC BY 4.0)}
\begin{document}
\maketitle

\noindent\emph{Source and licence.} Incompleteness in Quantified Conditional Logic. Version arXiv:2602.04073v1 (CC BY 4.0), distributed under the Creative Commons Attribution 4.0 International licence, as recorded in the arXiv licence metadata for this exact version and shown on the abstract page https://arxiv.org/abs/2602.04073v1. Full rights notice: licenses/arxiv-2602.04073-CC-BY-4.0.txt.\par\smallskip

\noindent\textit{Editorial scope.} Counted unit: the Claim whose source label is \texttt{boundary} in the article (source0.tex lines 1161--1169, source label \texttt{boundary}), delivered together with the article's own complete original proof of it (source0.tex lines 1171--1209, one \texttt{proof} environment of 39 lines). No mathematics is written, repaired or re-derived: statement and proof are the article's own text line for line. Required context retained uncounted: into-dnf (lines 1115--1124); observation (lines 1153--1155). Every definition, display and label of those lines is retained unchanged. Omitted from this package and not redistributed: 1--1114, 1125--1152, 1156--1160, 1170--1170, 1210--1503. Nothing inside a retained excerpt is omitted or altered: every retained body is delivered exactly as the article's lines are written, so no omission receipt of any kind accompanies this package. Citations: the retained excerpts carry no citation command at all, so no bibliography entry of the article is transcribed and no reference is redistributed. Reference adaptation: none was needed and none was made; every reference inside the delivered unit resolves inside it, so no unresolved reference or citation is produced. Printed slips kept literal: no printed word, symbol, hypothesis or number of the counted statement or of its proof was altered. Layout and class: the article's own class is \texttt{article} with packages including geometry, amsmath, amssymb, amsthm, enumitem, pifont, csquotes, tikz, caption, graphicx, biblatex; the wrapper is the lane's pinned \texttt{amsart} editorial wrapper at 10pt on A4, which restates the theorem-like environments the retained text needs and adds the article's own macro definitions that the retained text needs (\texttt{\textbackslash then}), with extra wrapper packages (bm, bbm, dsfont, xspace, cancel, mathtools). The delivered numbering is the unit's own. Assets: no retained span contains a figure environment, a graphics-inclusion command, a PDF-page inclusion, a table or a TikZ picture; no image file is copied into this package, the package declares no asset dependency and no image file is redistributed with it. Licence: version arXiv:2602.04073v1 of this article is distributed under the Creative Commons Attribution 4.0 International licence, as recorded in the arXiv licence metadata for this exact version and shown on the abstract page https://arxiv.org/abs/2602.04073v1. A full rights notice is delivered at licenses/arxiv-2602.04073-CC-BY-4.0.txt. Characters: every retained excerpt is pure ASCII, so the delivered engine, XeLaTeX with its default Latin Modern text font, reports no missing character.
\begin{claim}\label{into-dnf}
    Suppose that $\mathcal{K},-\infty,g \nVdash \phi$. Then $\phi$ is $\mathcal{K}$-equivalent to a disjunction in which each disjunct has of one of the following forms ($n,m \geq 0$):
    \begin{align*}
        & F(x_1) \wedge \cdots \wedge F(x_n) \wedge \neg F(y_1) \wedge \cdots \wedge \neg F(y_m) \\
        & F(x_1) \wedge \cdots \wedge F(x_n) \wedge \exists x \neg F(x) \\
        & \neg F(x_1) \wedge \cdots \wedge \neg F(x_n) \wedge \exists x F(x) \\
        & \exists x F(x) \wedge \exists x \neg F(x) \\
        & \forall x F(x)
    \end{align*}
\end{claim}

\begin{claim}\label{observation}
    If $m\geq n$ then $\mathcal{K} \Vdash F(m)\then F(n)$.
\end{claim}

\begin{claim}\label{boundary}
    Suppose that $\mathcal{K},-\infty,g \nVdash \phi$. Then $\phi$ is $\mathcal{K}$-equivalent to a disjunction in which each disjunct has one of the following forms:
    \begin{align*}
        & F(j) \wedge \neg F(j+1) \text{ where } k \leq j <-1 \\
        % & \exists x F(x) \wedge \exists x \neg F(x) \\
        & \forall x F(x) \\
        & \neg F(k) \wedge \exists x F(x)
    \end{align*}
\end{claim}

\begin{proof}
    We assume that $\phi$ is in the form provided by Claim~\ref{into-dnf}. 
    We now consider the different disjuncts of $\phi$ case by case. 
    Note that $\forall x F(x)$ is already in the appropriate form. 
    We now consider the other cases.
    
    \textbf{Case 1:} Consider some instance of:
    $$(B) \quad F(x_1)\wedge\dots\wedge F(x_n) \wedge \neg F(y_1) \wedge \dots \wedge \neg F(y_m)$$
    Without loss of generality, let's stipulate that $g(x_1)$ is greatest among the $g(x_i)$s and $g(y_1)$ is least among the $g(y_i)$s.
    
    As a first step, we note that $(B)$ is $\mathcal{K}$-equivalent to:
    $$(C) \quad F(x_1)\wedge \neg F(y_1)$$
    That $(B)$ $\mathcal{K}$-entails $(C)$ is obvious. 
    For the other direction, note first that $g(x_1)>g(x_i)$ for all $i\neq 1$. 
    Hence Claim~\ref{observation} implies that the conjunct $F(x_1)$ implies each $F(x_i)$. 
    Using Claim~\ref{observation} and the the fact $g(y_1)<g(y_i)$, we similarly see that the conjunct $\neg F(y_1)$ implies each $\neg F(y_i)$. 
    
    We now want to see that $(C)$ is $\mathcal{K}$-equivalent to a disjunction of claims of the form $F(j)\wedge\neg F(j+1)$. 
    Note that $(C)$ states that---relative to the variable assignment $g$---the boundary between $F$ and $\neg F$ occurs within the interval $[x_1,y_1]$. 
    Hence $(C)$ $\mathcal{K}$-entails the disjunction:
    $$(D) \quad \big(F(x_1)\wedge\neg F(x_1+1)\big)\vee \dots \vee \big(F(y_1-1)\wedge \neg F(y_1)\big)$$
    Appealing to Claim~\ref{observation} once again, we see that each disjunct of $(D)$ $\mathcal{K}$-entails $(C)$. 
    All of this is to say that $(B)$ is $\mathcal{K}$-equivalent to $(D)$.
    
    \textbf{Case 2:} Now a formula of the form 
    $$(E) \quad F(x_1)\wedge\dots\wedge F(x_n) \wedge \exists x \neg F(x)$$
    is $\mathcal{K}$-equivalent to a disjunction of claims of the form $F(j) \wedge \neg F(j+1)$. 
    For this we use the same reasoning as in the previous case. 
    We once again stipulate that $g(x_1)$ is greatest among the $g(x_i)$s. 
    Then, by applying Claim~\ref{observation} once again, we see that $(E)$ is $\mathcal{K}$-equivalent to the disjunction:
    $$\big(F(x_1)\wedge\neg F(x_1+1)\big)\vee \dots \vee \big(F(-2)\wedge \neg F(-1)\big)$$
    
    \textbf{Case 3:} Finally, let's consider formulas of the form:
    $$(F) \quad \neg F(x_1)\wedge\dots\wedge \neg F(x_n) \wedge \exists x F(x)$$
    Without loss of generality, let's stipulate that $g(x_1)$ is least among the $g(x_i)$s. 
    Then Claim~\ref{observation} immediately entails that $(F)$ is $\mathcal{K}$-equivalent to:
    $$\neg F(x_1)\wedge \exists x F(x)$$
    This completes the proof.
\end{proof}

\end{document}
